% ============================================================
%  QSS20 Final Project — Digital Health Exclusion in Kenya
%  Single-column PNAS mathematics template (pnasmathematics).
%  Start from the PNAS single-column Overleaf template, replace
%  its main.tex with this file and its .bib with references.bib.
%  In pnas-new.cls set displaywatermark to false.
% ============================================================

\documentclass[9pt,twoside]{pnas-new}

\articletype{SOCIAL SCIENCES}

\templatetype{pnasmathematics}

\begin{document}

\title{Who Gets Left Behind? Digital Exclusion Among Kenyan Women and
the Reach of Maternal Health Interventions}

\author[a]{Aulanni Kidd}

\affil[a]{QSS20, Dartmouth College, Hanover, NH 03755}

\leadauthor{Kidd}

\significancestatement{Digital maternal health tools are expanding
rapidly across sub-Saharan Africa, yet they assume a baseline of
connectivity that many women do not have. Using nationally representative
data from the Kenya 2022 Demographic and Health Survey ($n = 32{,}156$),
we show that education is the single strongest predictor of internet use,
stronger than wealth or geography. Women with no education have odds
of internet use roughly 52.6 times lower than women with higher
education. A rural woman with primary education has less than a 5\%
predicted probability of using the internet in the poorest wealth
quintile. These findings imply that internet-dependent maternal health
interventions, as currently designed, are least reachable by the women
who need them most.}

\authorcontributions{A.K.\ designed the study, cleaned the data,
conducted all analyses, and wrote the manuscript.}

\authordeclaration{The author declares no competing interests.}

\equalauthors{}

\correspondingauthor{\textsuperscript{a}Email:
aulanni.kidd@dartmouth.edu}

\keywords{digital health $|$ maternal health $|$ Kenya $|$ digital
divide $|$ mHealth}

\begin{abstract}
Digital maternal health tools (SMS-based reminders, mHealth
applications, and online health portals) are expanding rapidly across
sub-Saharan Africa. Yet these tools assume a baseline of digital
connectivity that many women do not have. Using data from the Kenya 2022
Demographic and Health Survey ($n = 32{,}156$ women aged 15--49), this
paper examines how wealth, education, geographical residence, and age
predict phone ownership and internet use. While 74.7\% of Kenyan women
own a phone, only 36.9\% use the internet, a gap that widens
dramatically across socioeconomic strata. Logistic regression reveals
that education is the single strongest predictor of internet use: women
with no education have odds of internet use roughly 52.6 times lower than
women with higher education, even after controlling for wealth, residence,
and age. Wealth is the second-largest predictor, with the
richest-quintile women nearly 20 times more likely to use the internet
than the poorest. Urban residence confers only a modest independent
advantage after controlling for socioeconomic composition. Predicted
probabilities show that a rural woman with primary education has less than
a 5\% chance of internet use in the poorest wealth quintile, rising to
only 48\% in the richest. These findings suggest that digital maternal
health interventions systematically exclude the women most at risk and
call for a phone-first, education-aware design philosophy.
\end{abstract}

\dates{This manuscript was compiled on \today.}
\doi{\url{www.pnas.org/cgi/doi/10.1073/pnas.XXXXXXXXXX}}

\maketitle
\thispagestyle{firststyle}
\ifthenelse{\boolean{shortarticle}}
  {\ifthenelse{\boolean{singlecolumn}}{\abscontentformatted}{\abscontent}}{}

\Firstpage

Maternal mortality remains one of the most persistent global health
inequities, with the vast majority of preventable deaths concentrated in
sub-Saharan Africa \cite{who2025,unicef2023}. In response, governments,
NGOs, and technology companies have invested heavily in digital health
tools designed to extend the reach of maternal care beyond traditional
clinical settings \cite{mwaisaka2025}. Kenya has been a particularly
active site of mHealth innovation: community health workers use mobile
devices to track pregnant patients, SMS systems remind women of antenatal
appointments, and mobile-linked platforms deliver nutrition and
complication-alert information \cite{mwaisaka2025,nyatuka2025}.

Yet a fundamental question underlies these interventions: do the women
they are designed to reach actually have access to the digital
infrastructure they require? If digital health tools disproportionately
reach women who are already wealthier, more educated, and more urban,
they risk reinforcing, rather than reducing, maternal health
inequities. Understanding the structure of digital exclusion is therefore
a precondition for designing interventions that work for the women who
need them most.

\Endparasplit

This paper makes two contributions. First, it provides a nationally
representative descriptive picture of phone ownership and internet use
among Kenyan women aged 15--49 across wealth, education, and residence
groups. Second, it uses logistic regression to identify the independent
predictors of digital access, disentangling the roles of wealth,
education, and geography. Results reveal that education is the dominant
gating factor for complex digital tools, that wealth drives an enormous
gradient in internet use, and that urban-rural disparities in
connectivity are largely explained by socioeconomic composition rather
than geography per se.

\section*{Background}

Digital health interventions in low- and middle-income countries have
shown promising results in controlled settings. Mwaisaka et al.\
\cite{mwaisaka2025} found that a combined mHealth and community-based
intervention in Kwale County, Kenya improved maternal outcomes for young
mothers. Nyatuka et al.\ \cite{nyatuka2025} argue for a frugal ICT4D
framework that prioritises low-cost, locally appropriate digital
solutions for maternal care. However, most published mHealth evaluations
are conducted in peri-urban settings and disproportionately enrol women
with secondary education or above
\cite{batist2019,till2023,soyege2025}, precisely the groups least
likely to represent Kenya's highest-risk maternal populations.

Prior literature has documented substantial digital divides in
sub-Saharan Africa by wealth, education, and geography, but few studies
use nationally representative samples to decompose these predictors
simultaneously. The Kenya 2022 DHS provides a rare opportunity to
quantify digital access at population scale, linking connectivity
variables to the full socioeconomic profile of respondents.

\section*{Data and Methods}

\subsection*{Data source}

This analysis uses the Kenya 2022 Demographic and Health Survey (DHS)
Individual Recode \cite{kdhs2022}, a nationally representative
cross-sectional household survey of women aged 15--49. The working sample
is $n = 32{,}156$ after listwise deletion of missing values across the
six analytic variables.

\subsection*{Variables}

Two binary outcomes are examined: \textit{phone ownership}
(\texttt{v169a}, owns a mobile telephone, 0/1) and \textit{internet use}
(\texttt{v171a}, uses the internet, 0/1). Phone ownership is the
hardware prerequisite for mobile-based interventions; internet use is
required for app- and web-based tools. Predictors are household
\textit{wealth quintile} (\texttt{v190}; 1~=~poorest, 5~=~richest),
\textit{education level} (\texttt{v106}; no education, primary,
secondary, higher), \textit{urban/rural residence} (\texttt{v025}), and
\textit{age} in years (\texttt{v012}). Table~\ref{tab:sample} reports
the sample distribution across all predictor categories.

\begin{table}[h]
\centering
\caption{Sample characteristics. Kenya DHS 2022 women aged 15--49.
$N = 32{,}156$.}
\label{tab:sample}
\begin{tabular}{@{}lrcc@{}}
\toprule
Group & $n$ & Phone & Internet \\
\midrule
\textit{Wealth quintile} & & & \\
\quad Q1 (Poorest) & 7{,}073 & 54.3\% &  5.4\% \\
\quad Q2            & 5{,}742 & 66.8\% & 17.1\% \\
\quad Q3            & 6{,}345 & 76.4\% & 31.9\% \\
\quad Q4            & 7{,}160 & 85.4\% & 53.7\% \\
\quad Q5 (Richest)  & 5{,}836 & 92.5\% & 79.5\% \\
\addlinespace
\textit{Education level} & & & \\
\quad No education  &  3{,}836 & 66.2\% &  4.9\% \\
\quad Primary       & 11{,}807 & 69.3\% & 16.8\% \\
\quad Secondary     & 11{,}634 & 73.0\% & 45.2\% \\
\quad Higher        &  4{,}879 & 98.8\% & 90.9\% \\
\addlinespace
\textit{Residence} & & & \\
\quad Rural         & 19{,}770 & 68.4\% & 24.3\% \\
\quad Urban         & 12{,}386 & 85.0\% & 56.9\% \\
\addlinespace
\textit{Full sample} & 32{,}156 & 74.7\% & 36.9\% \\
\bottomrule
\end{tabular}
\smallskip\\
\raggedright{\footnotesize Mean age = 28.6 years (SD not reported;
computed from \texttt{v012}).}
\end{table}

\subsection*{Analytic strategy}

The analysis proceeds in three steps. First, group means of phone
ownership and internet use are computed across all predictor categories.
Second, two logistic regression models, one per outcome, regress each
binary outcome on wealth quintile, education, residence, and age
simultaneously, with reference categories of wealth Q1, higher
education, and rural residence. Results are reported as odds ratios (ORs)
with 95\% confidence intervals. Third, predicted probabilities are
computed across wealth quintiles holding education fixed at Primary,
residence at Rural, and age at the sample mean (28.6 years). All
analyses were conducted in Python using \texttt{statsmodels}.

\section*{Results}

\subsection*{Overall access rates}

Across the full sample, 74.7\% of women own a mobile phone while only
36.9\% use the internet. This 37.8 percentage-point gap indicates that
phone ownership substantially overstates effective digital connectivity.

\subsection*{Descriptive patterns}

\textbf{Wealth.} Fig.~\ref{fig:wealth} shows that internet use rises
steeply from 5.4\% in Q1 to 79.5\% in Q5, a 74.1 percentage-point
gap. Phone ownership also rises with wealth (54.3\% to 92.5\%), but the
gradient is shallower (38.2~pp). In Q1, more than half of women own a
phone yet fewer than one in eighteen uses the internet.

\begin{figure}[H]
\centering
\includegraphics[width=0.7\linewidth]{figures/phone_vs_internet_wealth.png}
\caption{Phone ownership and internet use by wealth quintile. Q1 =
poorest, Q5 = richest. Kenya DHS 2022 ($n=32{,}156$).}
\label{fig:wealth}
\end{figure}

\textbf{Education.} Fig.~\ref{fig:education} shows the steepest gradient
of any predictor for internet use: from 4.9\% among women with no
education to 90.9\% among those with higher education, an 86.0~pp
gap. For phone ownership the gap is 32.6~pp (66.2\% to 98.8\%). The jump
from secondary to higher education is especially large for internet use
(45.2\% to 90.9\%), suggesting qualitatively different digital engagement
at the post-secondary level.

\begin{figure}[H]
\centering
\includegraphics[width=0.7\linewidth]{figures/phone_vs_internet_education.png}
\caption{Phone ownership and internet use by highest education level.
Kenya DHS 2022 ($n=32{,}156$).}
\label{fig:education}
\end{figure}

\textbf{Urban/rural residence.} Internet use is 56.9\% among urban women
versus 24.3\% among rural women, a 32.6~pp gap. Phone ownership shows
a smaller gap (85.0\% vs.\ 68.4\%), and 68.4\% of rural women own a
phone, indicating a viable foundation for SMS-based outreach.

\textbf{Compound disadvantage.} Fig.~\ref{fig:heatmap} shows internet
use rates at each wealth-by-education combination. A Q1 woman with no
education has a 1.5\% internet use rate; a Q5 woman with higher education
has 96.8\%. Education can partially offset low wealth: a Q1 woman with
higher education uses the internet at 54.8\%, far above the Q3 average
for women with no education (9.8\%). Conversely, wealth alone cannot
overcome an education deficit: Q5 women with no education use the
internet at only 32.4\%.

\begin{figure}[H]
\centering
\includegraphics[width=0.7\linewidth]{figures/heatmap_wealth_education.png}
\caption{Internet use rate (\%) by wealth quintile (rows) and education
level (columns). Kenya DHS 2022 ($n=32{,}156$).}
\label{fig:heatmap}
\end{figure}

\subsection*{Logistic regression}

Table~\ref{tab:regression} reports odds ratios from both models.
Fig.~\ref{fig:forest} visualises these on a log scale.

\begin{table*}[t]
\centering
\caption{Logistic regression odds ratios for phone ownership and internet
use. Reference categories: Wealth Q1, Higher education, Rural residence.
Kenya DHS 2022 ($n=32{,}156$). All ORs significant at $p<0.05$ except
age in the internet use model ($p=0.293$).}
\label{tab:regression}
\begin{tabular}{@{}lcccc@{}}
\toprule
& \multicolumn{2}{c}{\textbf{Internet Use}}
& \multicolumn{2}{c}{\textbf{Phone Ownership}} \\
\cmidrule(lr){2-3}\cmidrule(lr){4-5}
Predictor & OR & 95\% CI & OR & 95\% CI \\
\midrule
\textit{Wealth (ref: Q1)} & & & & \\
\quad Q2 & 2.221 & (1.947,\ 2.533) & 1.617 & (1.480,\ 1.766) \\
\quad Q3 & 4.305 & (3.800,\ 4.878) & 2.356 & (2.147,\ 2.585) \\
\quad Q4 & 8.673 & (7.617,\ 9.875) & 3.619 & (3.240,\ 4.041) \\
\quad Q5 & 19.894 & (17.176,\ 23.042) & 5.546 & (4.792,\ 6.419) \\
\addlinespace
\textit{Education (ref: Higher)} & & & & \\
\quad Secondary    & 0.132 & (0.118,\ 0.147) & 0.078 & (0.060,\ 0.101) \\
\quad Primary      & 0.044 & (0.039,\ 0.049) & 0.044 & (0.033,\ 0.057) \\
\quad No education & 0.019 & (0.016,\ 0.023) & 0.043 & (0.032,\ 0.057) \\
\addlinespace
\textit{Residence (ref: Rural)} & & & & \\
\quad Urban & 1.103 & (1.022,\ 1.190) & 1.161 & (1.067,\ 1.262) \\
\addlinespace
Age (per year) & 0.998 & (0.995,\ 1.002) & 1.138 & (1.133,\ 1.142) \\
\midrule
Pseudo $R^2$ & \multicolumn{2}{c}{0.357} & \multicolumn{2}{c}{0.273} \\
$N$ & \multicolumn{4}{c}{32{,}156} \\
\bottomrule
\end{tabular}
\end{table*}

\begin{figure}[H]
\centering
\includegraphics[width=0.7\linewidth]{figures/odds_ratio_plot.png}
\caption{Forest plot of odds ratios (log scale) for internet use and
phone ownership. Dashed line at OR~=~1. Reference: Wealth Q1, Higher
education, Rural. Kenya DHS 2022 ($n=32{,}156$).}
\label{fig:forest}
\end{figure}

\textbf{Education is the dominant predictor.} Women with no education
have odds of internet use roughly 52.6 times lower than those with higher
education (OR~=~0.019; 95\%~CI: 0.016--0.023), exceeding the largest
wealth contrast (Q5 vs.\ Q1: OR~=~19.894). The education effect is
larger for internet use than phone ownership (OR~=~0.043 for phones),
indicating that education shapes the \textit{ability to use} digital
services beyond the ability to acquire a device.

\textbf{Wealth drives a massive gradient.} Controlling for education and
residence, richest-quintile women have nearly 20 times the odds of
internet use versus poorest-quintile women. The gradient is monotonic and
accelerating across quintiles (Q2:~2.22; Q3:~4.31; Q4:~8.67;
Q5:~19.89). The wealth effect on phone ownership is significant but
substantially smaller (Q5 vs.\ Q1: OR~=~5.55).

\textbf{Urban residence confers modest independent advantage.} After
controlling for wealth and education, urban women have only 10.3\% higher
odds of internet use than rural women (OR~=~1.103; $p=0.012$). This
small effect, relative to the large raw urban-rural gap (32.6~pp),
implies that most of the observed urban advantage is compositional,
driven by urban areas concentrating wealthier, better-educated
populations, rather than reflecting infrastructure differences per se.

\textbf{Age effects diverge across outcomes.} Age has a significant
positive effect on phone ownership (OR~=~1.138 per year; $p<0.001$),
consistent with longer lifetime exposure to mobile adoption. By contrast,
age has a near-zero and non-significant effect on internet use
(OR~=~0.998; $p=0.293$), indicating that internet adoption is driven by
wealth and education rather than generational patterns.

\subsection*{Predicted probabilities}

Table~\ref{tab:predprobs} shows predicted probabilities across wealth
quintiles for a rural woman with primary education at the sample mean age,
a profile representative of a large share of Kenya's highest-risk
maternal population. Even the wealthiest rural woman with only primary
education has a predicted internet use probability below 50\% (48.2\%),
while the poorest has just 4.5\%. Phone ownership remains high across
the distribution (54.8\%--87.0\%), underscoring the phone-internet gap.

\begin{table}[H]
\centering
\caption{Predicted probabilities by wealth quintile. Profile: Rural,
Primary education, Mean age (28.6 years). Kenya DHS 2022.}
\label{tab:predprobs}
\begin{tabular}{@{}lccc@{}}
\toprule
Quintile & Phone & Internet & Gap \\
\midrule
Q1 (Poorest) & 54.8\% & 4.5\%  & 50.3 pp \\
Q2           & 66.2\% & 9.4\%  & 56.8 pp \\
Q3           & 74.0\% & 16.8\% & 57.2 pp \\
Q4           & 81.4\% & 28.9\% & 52.5 pp \\
Q5 (Richest) & 87.0\% & 48.2\% & 38.8 pp \\
\bottomrule
\end{tabular}
\end{table}

\section*{Discussion}

\textbf{Education as the primary bottleneck.} The finding that education
dominates wealth as an independent predictor of internet use has direct
implications for digital health design. Many connectivity interventions
assume that reducing the cost of devices or data will unlock access for
underserved women. This analysis suggests that for women with low
education, cost is not the binding constraint: digital literacy (the
ability to navigate applications and interpret health information online)
may be a prerequisite that cost-based interventions alone cannot supply.

\textbf{Phone ownership as a viable foundation.} Despite the connectivity
gap documented here, 68.4\% of rural women own a phone, and phone
ownership is relatively stable across much of the wealth distribution.
This points to a viable infrastructure for reaching marginalised women
through SMS-based and voice-based interventions that do not require
internet access, as demonstrated by Mwaisaka et al.\ \cite{mwaisaka2025}.
The findings lend quantitative support to a phone-first design philosophy
for maternal health outreach.

\textbf{Compositional vs.\ infrastructural urban advantage.} The finding
that the urban-rural gap largely disappears after controlling for wealth
and education reframes the policy question. It suggests that equalising
rural women's educational attainment and economic resources, rather than
simply expanding rural network coverage, may be the more effective lever
for closing the connectivity gap.

\textbf{Limitations.} This analysis is cross-sectional and cannot
establish causality. Outcome variables are binary self-reports and do not
capture frequency or type of digital use. Standard errors do not account
for the DHS cluster sampling design and may be slightly underestimated.
Subgroup heterogeneity by region or age cohort is not explored.

\section*{Conclusion}

Digital maternal health tools carry the promise of extending quality care
to Kenya's most vulnerable women. This analysis demonstrates that the
women who face the highest risks of maternal mortality (those who are
poorest, least educated, and most rural) are precisely those with the
lowest rates of internet use. Education is the single largest barrier;
wealth is the second. Urban-rural gaps are largely explained by
socioeconomic composition. For internet-dependent tools to achieve
equity, they must either be paired with intensive digital literacy
support or radically simplified to reduce the education threshold
required. Without deliberate design for digital exclusion, maternal
health interventions risk reaching those already advantaged, leaving
behind the women who need them most.

\dataavail{All analysis code is available at
\url{https://github.com/aulanni-kidd/qss20_final_digital_health_access}.
The underlying data are the Kenya 2022 DHS Individual Recode, available
from the DHS Program (\url{https://dhsprogram.com}) upon registration.}

\acknow{Data were obtained from the Kenya 2022 Demographic and Health
Survey, conducted by the Kenya National Bureau of Statistics and ICF.
This project was completed as part of QSS20 at Dartmouth College, Spring
2026.}

\showacknow{}

\bibliography{references}

\end{document}